% Source: https://artofproblemsolving.com/community/c3h1762366_mathcounts_handbook
% Stable identity: aops-topic-1762366-q148
In the order $W,X,Y,Z$, the vertices $W$ and $Y$ are opposite, so $WY$ is a diagonal. (They need not be adjacent, which is why ``$W$ and $Y$ do not form a side'' is no obstacle.) The centre is the midpoint $M=(3,4)$, and $\overrightarrow{MW}=(-1,-3)$. The other two vertices are obtained by rotating this vector through $\pm90^\circ$: $(3,-1)$ and $(-3,1)$. So $X$ and $Z$ are $(6,3)$ and $(0,5)$ in some order.

\begin{asy}
usepackage("amsmath");
size(6.5cm);
defaultpen(fontsize(10pt));
import graph;
pair Wp=(2,1),X=(6,3),Y=(4,7),Z=(0,5),M=(3,4);
draw((-0.8,0)--(7.2,0),Arrow);draw((0,-0.5)--(0,8),Arrow);
label("$x$",(7.2,0),S);label("$y$",(0,8),W);
draw(Wp--X--Y--Z--cycle,linewidth(0.9));draw(Wp--Y,dashed);draw(X--Z,dashed);
dot(Wp);dot(X);dot(Y);dot(Z);dot(M);
label("$W(2,1)$",Wp,dir(-40));label("$X(6,3)$",X,E);label("$Y(4,7)$",Y,N);label("$Z(0,5)$",Z,W);label("$(3,4)$",M,1.5*dir(27));
label("\mbox{}",(min(currentpicture,user=true)+max(currentpicture,user=true))/2,black); // keeps the text after the figure black when LaTeX typesets the labels inline
\end{asy}

Either way $|ab-cd|=|6\cdot3-0\cdot5|=\boxed{18}$.
